\documentclass[11pt,a4paper]{article}
\usepackage[margin=1.8cm]{geometry}
\usepackage{amsmath,amssymb,tabularx,array,needspace,fontspec,enumitem,cancel}
\newcommand{\tb}{\cancel{||||}\,}
\setmainfont{FreeSerif}
\newfontfamily\urdufont[Script=Arabic]{FreeSerif}
\TeXXeTstate=1
\newcommand{\urdu}[1]{{\urdufont\beginR #1\endR}}
\setlength{\parindent}{0pt}
\setlength{\parskip}{0pt}
\renewcommand{\arraystretch}{1.15}
\pagestyle{plain}

% \mcq{number}{question}{A}{B}{C}{D}
\newcommand{\mcq}[6]{%
  \Needspace{4\baselineskip}%
  \par\medskip\noindent\textbf{#1.}\ #2\par\nopagebreak
  \noindent\begin{tabularx}{\linewidth}{@{\hspace{1.5em}}*{4}{>{\raggedright\arraybackslash}X}@{}}
    (A)~#3 & (B)~#4 & (C)~#5 & (D)~#6
  \end{tabularx}\par}

\begin{document}

\begin{center}
  {\Huge\bfseries The Hope Science Academy}\\[3pt]
  {\small Excellence is our standard not goal}\\[2pt]
  {\small 247-E-1, (Opposite LDA School) Johar Town, Lahore. PH: 0300-8194789, 0322-4157001}
\end{center}

\medskip
\noindent
\begin{tabularx}{\linewidth}{|X|X|X|X|}
\hline
\textbf{Student Name:} & \textbf{Roll Num:} & \textbf{Class:} 12th & \textbf{Paper Code:} \\ \hline
\textbf{Subject Name:} Business Statistics & \textbf{Time Allowed:} & \textbf{Total Marks:} 347 & \textbf{Exam Date:} 02-Nov-2026 \\ \hline
\multicolumn{4}{|l|}{\textbf{Exam Syllabus:} FULL BOOK} \\ \hline
\end{tabularx}

\bigskip
\noindent\textbf{Q1. Choose the correct answer.}\hfill\textbf{1$\times$347=347}

\mcq{1}{The word statistics may have been derived from the gGerman word:}{Statistics}{Statistic}{Status}{Statistik}
\mcq{2}{The word statisics come from Latin word:}{Status}{Statistics}{Statista}{Statistique}
\mcq{3}{The word statistics is taken from:}{Latin word}{Greek word}{French word}{English word}
\mcq{4}{Statistics is a branch of:}{Applied economics}{Applied mathematics}{Applied philosophy}{None of these}
\mcq{5}{Statistics is a backbone of:}{Mathematics}{Computer}{Research}{Accounting}
\mcq{6}{The word statistics is now a day used in:}{One way}{Two ways}{Three ways}{None of these}
\mcq{7}{The word STATISTICS was first used by scholar:}{Laspeyre}{Bowley}{Fisher}{Gott achenwall}
\mcq{8}{In plular senee, statistics mean:}{Aggregate of fact}{Methodology}{Meteorology}{None of these}
\mcq{9}{Statistics are always:}{Aggregate of facts}{Individual}{Fixed}{None}
\mcq{10}{Research organization provide:}{Primary data}{Secondary data}{Constant}{Continuous data}
\mcq{11}{Statistics lawas are valid on:}{Long run}{Short run}{Medium run}{None of these}
\mcq{12}{Brach of statistics which deals with the collection, presentation and analysis of data is called:}{Descriptive statistics}{Statics}{Inferential statistics}{Statistics}
\mcq{13}{The total group under discussion:}{Sample}{Population}{Subset}{Variable}
\mcq{14}{Part of population under study is called:}{Variable}{Data}{Sample}{Parameter}
\mcq{15}{Representative part of population is called:}{Sample}{Population}{Total group}{None of these}
\mcq{16}{Pie is a:}{Constant}{Variable}{Statistics}{Coefficient}
\mcq{17}{A constant can assume:}{Only one value}{More than one value}{Different value}{No value}
\mcq{18}{The characteristics which varies from one object to another object is called:}{Variable}{Constant}{Co-efficient}{Exponent}
\mcq{19}{A collection of some elements from a group is called:}{Census}{Population}{Sample}{Registration}
\mcq{20}{U is called:}{Variable}{Statistics}{Parameter}{Coefficient}
\mcq{21}{Population census is conducted through:}{Sample survey}{Complete enumeration}{Accruing}{Investigation}
\mcq{22}{Statistics must be:}{Comparable}{Not comparable}{Discrete nature}{Qualitative nature}
\mcq{23}{Data in the census report are:}{Primary data}{Secondary data}{Arrayed data}{None of these}
\mcq{24}{Height of a student in a class is:}{Discrete variable}{Continuous variable}{Qualitative variable}{None of these}
\mcq{25}{The measure which is calculated from sample is called:}{Statistic}{Statistics}{Constant}{Variable}
\mcq{26}{Which of the following is example of discrete variable:}{Number of days in a week}{Heights of students}{Shoes size of college student}{None of these}
\mcq{27}{A qualitative variable is called:}{Parameter}{Statistic}{Attribute}{None of these}
\mcq{28}{Measurement usually provides:}{Discrete data}{Continuous data}{Qualitative}{Primary data}
\mcq{29}{A variable which can take any value between two numbers is called:}{Discrete variable}{Continuous variable}{Constant}{None of these}
\mcq{30}{Parameters are related to:}{Sample}{Population}{Mean}{None of these}
\mcq{31}{Weighted of whole earth is:}{Discrete variable}{Qualitative variable}{Constant}{None of these}
\mcq{32}{Secondary data cam be collected by:}{Registration}{Official source}{Direct personal observation}{None of these}
\mcq{33}{Colour of follower is an example of:}{Continuous variable}{Quantitative variable}{Qualitative variable}{None of these}
\mcq{34}{A value which does not change is called:}{Variable}{Statistic}{Constant}{Probability}
\mcq{35}{The sample means X is a:}{Parameter}{Stattistic}{Population}{Sample}
\mcq{36}{The pupose of the sample is to draw inference about:}{Sample}{Population}{Individual}{Art}
\mcq{37}{A measurable quantity is:}{Continuous}{Discrete}{Qualitative}{All of these}
\mcq{38}{Numerical value calculated from population is called:}{Parameter}{Statistics}{Statistic}{Constant}
\mcq{39}{Numbers of accidents is a:}{Discrete data}{Qualitative data}{Continuous data}{None of these}
\mcq{40}{The number of chairs is an example of:}{Constant}{Continuous variable}{Discrete variable}{Both a and b}
\mcq{41}{The life time of TV tube is:}{Discrete data}{Continuous data}{Constant}{None of these}
\mcq{42}{Data classified by attributes is called:}{Quantitative data}{Qualitative data}{Continuous data}{Discrete data}
\mcq{43}{A single numerical fact is called:}{Statistocs}{Constant}{Datum}{Data}
\mcq{44}{Registering a marriage by Nikah registrar is a source of:}{Primary data}{Secondary data}{Private data}{Fictitious data}
\mcq{45}{Personal investigation is a source of:}{Primary data}{Secondary data}{Discrete data}{Qualitative data}
\mcq{46}{Investigation through questionnaire collects:}{Primary data}{Secondary data}{Group data}{Qualitative data}
\mcq{47}{Questionaire method is used in collection:}{Primary data}{Secondary data}{Published data}{Both a and b}
\mcq{48}{Blood colour of a student is:}{Discrete variable}{Attribute}{Constant}{Both A and C}
\mcq{49}{The data which have undergone any treatment previously is called:}{Primary data}{Secondary data}{Continuous data}{None of these}
\mcq{50}{Weight of student in a class make:}{Discrete data}{Constant}{Qualitative data}{Continuous data}
\mcq{51}{The data collected from an individual is called:}{Secondary data}{Primary data}{Group data}{None of these}
\mcq{52}{The source of primary data are:}{Surveys}{Areas}{Unites}{Factors}
\mcq{53}{The data which have not under gone any statistical treatment are:}{Primary data}{Secondary data}{Discrete data}{Continuous data}
\mcq{54}{The group data is also called:}{Raw data}{Primary data}{Secondary data}{Qualitative data}
\mcq{55}{Printed data are always:}{Primary data}{Secondary data}{Confidential data}{Fictitious data}
\mcq{56}{Information recorded in original form are called:}{Continuous data}{Raw data}{Discrete data}{Arrayed data}
\mcq{57}{The data obtained from college record is:}{Primary}{Raw}{Secondary}{Qualitative}
\mcq{58}{The data which is not fresh in the characteristics is called:}{Ungrouped data}{Secondary data}{Primary data}{None of these}
\mcq{59}{Milk produced by a cow is the example of:}{Qualitative variable}{Discrete variable}{Continuous variable}{Constant}
\mcq{60}{The number od vowels in each word of this book are:}{Discrete data}{Continuous data}{Qualitative data}{Infinite data}
\mcq{61}{Number of flower on a tree is varibale:}{Continuous}{Qualitative}{Discrete}{None of these}
\mcq{62}{The statistics which is not measurable is called:}{An attribute}{Discrete}{Continuous}{None of these}
\mcq{63}{Data obtained from a newpaper is:}{Primary data}{Secondary data}{Raw data}{Extreme data}
\mcq{64}{Number of week in a month is an example of:}{Continuous variable}{Discrete variable}{Attribute}{None of these}
\mcq{65}{Issuing a national identity card is a:}{Registration}{Census}{Sampling}{Investigation}
\mcq{66}{The data published by economic survey of Pakistan are:}{Confidential}{Fictitious}{Primary}{Secondary}
\mcq{67}{The data which has already been collected by some one is called:}{Raw data}{Secondary data}{Array data}{Fictitious data}
\mcq{68}{Uniform of a student is an example of:}{Constant}{Variable}{Parameter}{Statistics}
\mcq{69}{Result declared are always:}{Primary data}{Secondary data}{Rowdata}{Extreme data}
\mcq{70}{Data on passports are:}{Primary}{Secondary}{Un-official}{Fictitious}
\mcq{71}{Size of the population as compare to sample is:}{Small}{Medium}{Large}{None of these}
\mcq{72}{The data obtained by conducting a survey is called:}{Primary data}{Secondary data}{Continuous data}{Qualitative data}
\mcq{73}{Number of student in a class is an example of:}{Class interval}{Qualitative data}{Continuous variable}{Discrete variable}
\mcq{74}{Classification in the process of arranging data according to:}{Two characteristics}{Differentcharacteristics}{Similarcharacteristics}{None of these}
\mcq{75}{When data is arranged in regular intervals of time, the classification is:}{Quantitative}{Quantitative}{Geographical}{Chronological}
\mcq{76}{Important bases of classification are:}{Two}{Three}{Four}{None of these}
\mcq{77}{A statistical table has at least:}{Four parts}{Two parts}{Five parts}{All of these}
\mcq{78}{Title should be in:}{Small letters}{Capital letters}{Italic letters}{Roman letter}
\mcq{79}{The part that contain row caption:}{Box head}{Stub}{Foot notes}{Body of the table}
\mcq{80}{Columncaption of a Table is known as:}{Stub}{Box head}{Foot note}{Body}
\mcq{81}{A source note and foot note are the:}{Same}{Different}{Complex}{Easy}
\mcq{82}{The portion of table containing column caption is called:}{Stub}{Box head}{Foot note}{Prefatory note}
\mcq{83}{Row caption is also called:}{Body}{Box head}{Stub}{Title}
\mcq{84}{Data grouped according to their magnitude is called:}{Sampling distriution}{Graph}{Frequency distribution}{Ogive}
\mcq{85}{The process of arranging data into rows and column is called:}{Frequency distribution}{Classification}{Tabulation}{Array}
\mcq{86}{The difference between two consecutive upper class limits is:}{Class marks}{Class frequency}{Class limits}{Class interval}
\mcq{87}{Mid-point of class limit is also called:}{Class mark}{Class limits}{Class boundaries}{None of these}
\mcq{88}{The numbers of observation falling in a particular class is called:}{Class mark}{Class interval}{Mid point}{None of these}
\mcq{89}{The size of a class is also called:}{Class interval}{Class frequency}{Class mark}{Class boundary}
\mcq{90}{The difference between upper and lower class boundaries of a class is called:}{Class interval}{Class limit}{Class marks}{Class frequency}
\mcq{91}{If there is no gap between consecutive classes, the limits are called:}{Class limits}{Class boundaries}{Class interval}{Class marks}
\mcq{92}{The mid point of the group 5.5 - 7.5 is:}{2}{5.5}{6}{6.5}
\mcq{93}{The lower and upper class limits are 20 and 30, the mid point of the class is:}{50}{30}{25}{20}
\mcq{94}{Main classes/ types of the graphs are:}{One}{Two}{Three}{Four}
\mcq{95}{The frequency of a class divided by the total frequency is called:}{Class frequency}{Relative frequency}{Total frequency}{Cumulative frequency}
\mcq{96}{Relative frequency can never be:}{Less than 1}{Equal to one}{0}{None of these}
\mcq{97}{Sum of all the relative frequencies is equal to:}{0.5}{100}{-1}{1}
\mcq{98}{In classes 15-19,20,24, the size of class interval is:}{3}{4}{5}{6}
\mcq{99}{The graph of frequency distribution is called:}{Histogram}{Historigram}{Ogive}{None of these}
\mcq{100}{While drawing a histogram, area of each bar is directly proportional to the:}{Mid point}{Frequency}{Total frequency}{Cumulative frequency}
\mcq{101}{In the class 60-70, class marks is:}{60}{10}{68}{65}
\mcq{102}{The graphs of cumulative frequency distribution is called:}{Histogram}{Ogive}{Frequency polygon}{Multiple bar chart}
\mcq{103}{Graph of time series is called:}{Ogive}{Historigram}{Histogram}{None of these}
\mcq{104}{A frequency polygon is closed figure of:}{One side}{Two side}{Three side}{Many side}
\mcq{105}{Both length and width are used in:}{One dimensional diagram}{Twodimensional diagram}{Threedimensional diagram}{Multidimensional diagram}
\mcq{106}{The budgets of two familes can be compared by:}{Subdivided rectangles}{Pie diagram}{Both A \& B}{Histogram}
\mcq{107}{An ogive is a:}{Cumulative frequency polygon}{Frequency polygon}{Frequency curve}{Histogram}
\mcq{108}{A pie diagram is represented by a:}{Circle}{Suqare}{Triangle}{Rectangle}
\mcq{109}{Angles are calculated in construction of:}{Bar cart}{Histogram}{Historigram}{Pie chart}
\mcq{110}{Total angle of pie chart are:}{$360^{\circ}$}{$180^{\circ}$}{$190^{\circ}$}{$90^{\circ}$}
\mcq{111}{A representative central value of the data:}{Average}{Dispersion}{Skewness}{Class interval}
\mcq{112}{Average is also called measures of:}{Variation}{Location}{Skewness}{Median}
\mcq{113}{It is a type of average:}{Mean}{Median}{Mode}{All of these}
\mcq{114}{Which average is highly affected by the extreme value:}{Mode}{A.M}{Median}{H.M}
\mcq{115}{Average of a constant data is:}{0}{The same constant}{Positive}{Negative}
\mcq{116}{Sum of value divided by their number is called:}{Mode}{Median}{Mean}{None of these}
\mcq{117}{Arithmetic mean can not be:}{Negative}{Positive}{Zero}{Imaginary}
\mcq{118}{The arithmetic mean is based on:}{All values}{Small values}{Extreme values}{None of these}
\mcq{119}{The mean of first n natural number is:}{$\dfrac{n(n+1)}{2}$}{$\dfrac{n}{2}$}{$\dfrac{n+1}{2}$}{$\dfrac{n-1}{2}$}
\mcq{120}{Most suitable average for skew Data is:}{Mean}{Median}{Mode}{None of these}
\mcq{121}{...... is based on all observation:}{A.M}{Median}{Mode}{None of these}
\mcq{122}{Arrangement of data is necessary in finding:}{Mode}{Mean}{Median}{All of these}
\mcq{123}{The mean of a constant ``a'' is equal to:}{0}{$\dfrac{a}{2}$}{a}{No possible}
\mcq{124}{...... average is affected by extreme values:}{A.M}{G.M}{H.M}{Median}
\mcq{125}{The notation for mean is:}{$\hat{X}$}{$\overline{X}$}{$\tilde{X}$}{$\hat{X}_{w}$}
\mcq{126}{Arithmetic means is effected by change of:}{Origin}{Scale}{Both}{None}
\mcq{127}{The sum of deviation is zero when deviation are takes from:}{Mean}{Median}{Weighted mean}{Mode}
\mcq{128}{Sum of deviations from aithmetic mean is always:}{Maximum}{Minimum}{Zero}{-1}
\mcq{129}{$\sum(X-X)=$ ................:}{Zero}{One}{Two}{None of these}
\mcq{130}{Is mean of 10 observations is 60 then sum of observations is:}{100}{60}{600}{50}
\mcq{131}{The sum of square deviations from mean is:}{Zero}{Minimum}{Negative}{None of these}
\mcq{132}{For an open and frequency distribution.It is not possible to calculate:}{Arithmetic mean}{Median}{Mode}{All of these}
\mcq{133}{If $\sum X=60$ and $\overline{X}=12$, than $N=?$ :}{4}{5}{6}{7}
\mcq{134}{If If $\sum(x-20)=25,\ \sum(x-18)=0$ then arithmetic meas is:}{20}{25}{18}{Zero}
\mcq{135}{Mean of symmetrical distribution is 90 value of its median will be:}{80}{85}{90}{75}
\mcq{136}{If $y=2x+5$ and $\overline{X}=5$ and $\overline{y}$ will be ''}{10}{12}{15}{20}
\mcq{137}{The sum of squared deviation is least from:}{Median}{Mean}{H.M}{G.M}
\mcq{138}{The mean of 5 values is 10 then sum will be:}{2}{15}{50}{$\dfrac{5}{10}$}
\mcq{139}{If mean of 5 observation is 40, then the sum of observations is:}{Zero}{200}{45}{35}
\mcq{140}{If $\sum(x-20)=0$ for 15 values, then the value of A.M is:}{15}{20}{25}{30}
\mcq{141}{Which one is equal to zero:}{$\sum(X-\overline{X})$}{$\sum(X-\overline{X})^{2}$}{$\sum(X-0)$}{$\sum(X-5)$}
\mcq{142}{If If $\sum_{i=1}^{10}(x_{i}-50)=100$, then $\overline{X}=$}{10}{50}{60}{100}
\mcq{143}{If $\overline{X}=10$ and $n=5$ then $\sum X=$ :}{15}{50}{2}{100}
\mcq{144}{If $\overline{X}=25$, which of the following is minimum:}{$\sum(x-27)^{2}$}{$\sum(x-25)^{2}$}{$\sum(x-22)^{2}$}{$\sum(x+25)^{2}$}
\mcq{145}{Change of origin and scale is used for calculation:}{Arithmetic mean}{Median}{Mode}{All of these}
\mcq{146}{If the arithmetic mean of two number $x_{1}$ and $x_{2}$ is 5, if $x_{1}=3$ then $x_{2}$ is:}{10}{7}{5}{30}
\mcq{147}{$\dfrac{\sum wx}{\sum w}$ is:}{A.M}{Mode}{Weighted Mean}{Median}
\mcq{148}{When the values are not of equal importance then we compute:}{Simple mean}{Mode}{Combined mean}{Weighted mean}
\mcq{149}{The value which divides arranged data into two equal halves:}{Mean}{Median}{Mode}{Probability}
\mcq{150}{Median is a value that divides an arrayed data into:}{5 equal parts}{2 equal parts}{10 equal parts}{None of these}
\mcq{151}{The notation for median is:}{$\overline{X}$}{$\tilde{X}$}{$\hat{X}$}{$\overline{X}_{w}$}
\mcq{152}{The central value of arranged data is called:}{Median}{Mean}{Mode}{None of these}
\mcq{153}{We must arrange data before calculating:}{A.M}{Median}{Mode}{None of these}
\mcq{154}{Median of the values -3,-2,0,1,2:}{0}{1}{2}{-2}
\mcq{155}{Which is a positional average:}{H.M}{Median}{A.M}{G.M}
\mcq{156}{The median for the data -2,0,2,5,-1,4 is:}{4}{1}{3.5}{Impossible}
\mcq{157}{If $Z=x-y$, the $\overline{Z}$.........:}{$\overline{X}+\overline{Y}$}{$\overline{X}-\overline{Y}$}{0}{$\overline{X}\ \overline{Y}$}
\mcq{158}{Median of series 5,2,3,4,2 is:}{4}{5}{3}{None of these}
\mcq{159}{Median of the data 2,3,4,5,6,7,8 is:}{2}{3}{4}{5}
\mcq{160}{The median of 3,4,5,6,9,10,12 is:}{5}{6}{9}{5.5}
\mcq{161}{The median of the data -2,0,2,5,-4,4 is:}{4}{1}{3.5}{Impossible}
\mcq{162}{What is the median of 6,8,10,12,14:}{13}{10}{16}{35}
\mcq{163}{Middle most value of the m observations arranged in ascending or descending order of magnitude is called:}{Mode}{Mean}{Median}{Combine Mean}
\mcq{164}{Cumulative frequency is calculated to find:}{A.M}{Median}{Mode}{G.M}
\mcq{165}{The most popular value of data is called:}{Mean}{Mode}{Median}{None}
\mcq{166}{A value which occurs greatest number of time in data:}{Mean}{Median}{Mode}{None of these}
\mcq{167}{Mode is not affected by:}{Smal values}{Large values}{Extreme values}{All values}
\mcq{168}{To find out the size of shoes, the best average is:}{Weighted mean}{Mean}{Mode}{None of these}
\mcq{169}{Mode of constant is:}{Zero}{Equal to that constant}{Undefined}{Not exist}
\mcq{170}{If a distribution has two modes then it is called:}{Bi model}{Uni model}{Tri model}{None of these}
\mcq{171}{If a distribution has one mode, then it is called:}{Uni-model}{Bi-model}{Tri model}{Multimodel}
\mcq{172}{Mode is a value which is:}{Central value}{Biggest value}{Most repeated value}{Smallest value}
\mcq{173}{Extreme scores on an examination have the following:}{They tend to rate it}{They tend to lower it}{They have no effect on it}{Difficultto tell}
\mcq{174}{The mode for the values 4,4,5,6,3,3,2 is:}{2}{4}{3}{3 and 4}
\mcq{175}{The mode of the values 2,2,2,0,0,0,1,1,7,9 is:}{2}{0}{9}{Both a \& b}
\mcq{176}{The most frequent value the data if exist is called:}{Arithmetic mean}{Median}{Mode}{None of these}
\mcq{177}{The model letter of the word STATISTICS is:}{S}{T}{Both S and T}{Both S and I}
\mcq{178}{The suitable average for qualitative data is:}{Mean}{Median}{Mode}{None of these}
\mcq{179}{Which of the following average have more than one value:}{A.M}{H.M}{Median}{Mode}
\mcq{180}{Mode of the series is 2,3,4,3,5,3,6,4:}{5}{4}{3}{6}
\mcq{181}{Mode of the data 3,6,9,3,4,5,2,3 is:}{4}{9}{20}{3}
\mcq{182}{In symmetrical distribution mean,median and mode are always:}{Unequal}{Equal}{Negative}{None of these}
\mcq{183}{If $\overline{X}=82$ and Median$=78$ then Mode is:}{50}{60}{70}{80}
\mcq{184}{Model letter of the word Population is:}{P}{N}{O}{Both a \& c}
\mcq{185}{The empirical relationship between mean, median and mode is, made=:}{3 mean - 2 median}{2 mean - 3 median}{3 median - 2 mean}{2 median - 3 mean}
\mcq{186}{index numbers are called the barometers of:}{Statistics}{Econimics}{Mathematics}{None of these}
\mcq{187}{By dividing price of current year by price of base year we get:}{Price relative}{Link relative}{Chain relative}{Close relative}
\mcq{188}{Relatives computed from fixed base are:}{Price relatives}{Link relatives}{Chain relatives}{None of these}
\mcq{189}{In fixed base method, the base price should be:}{Far-away}{Normal}{Abnormal}{None of these}
\mcq{190}{Index number for base year period is always:}{Mean}{Mode}{Unknown}{100}
\mcq{191}{$\dfrac{P_{n}}{P_{o}}\times100$ is equal to:}{Link relatives}{Chain relatives}{Price relatives}{None of these}
\mcq{192}{Index numbers are basically classified into:}{2 categories}{3 categories}{4 categories}{5 categories}
\mcq{193}{Index numbers are expressed in:}{Ratios}{Squares}{Combinations}{Percentages}
\mcq{194}{Price relatives are used in:}{Fixed base method}{Chain base method}{Fisher's method}{Consumer's price method}
\mcq{195}{Price relatives are equal to:}{$\dfrac{p_{n}}{p_{0}}\times100$}{$\dfrac{p_{n}}{p_{n-1}}\times100$}{$\dfrac{p_{o}}{p_{n}}\times100$}{$\dfrac{p_{n-1}}{p_{n}}\times100$}
\mcq{196}{Price index numbers calculate changes in:}{Price}{Quality}{Price \& quality}{None of these}
\mcq{197}{Price relatives $=P_{n}\times100$}{$P_{n}$}{$P_{o}$}{$P_{n-1}$}{$P_{n-2}$}
\mcq{198}{Price relatives and link relatives are:}{Same}{Different}{100}{Zero}
\mcq{199}{Link relatives are used in:}{Fixed base method}{Chain base method}{Both a \& b}{None of these}
\mcq{200}{Link relatives=:}{$\dfrac{P_{o}}{P_{n-1}}\times100$}{$\dfrac{P_{n}}{P_{n-1}}\times100$}{$\dfrac{P_{n}}{P_{n-2}}\times100$}{$\dfrac{P_{n}}{P_{o}}\times100$}
\mcq{201}{Chain base method is batter than fixed base:}{True}{False}{Partially true}{None of these}
\mcq{202}{The ratio obtained by dividing the price in a given year by the price in proceeding year are called:}{Volume relatives}{Price relatives}{Link relatives}{None}
\mcq{203}{In calculating link relatives, the base period is:}{Fixed}{Constant}{Not fixed}{Any year}
\mcq{204}{In chain base method, we use link-relatives:}{Always}{Sometimes}{Never}{None of these}
\mcq{205}{For computing chain index numbers, we compute:}{Price relatives}{Link relatives}{Weighted indicies}{Fixed base ondicies}
\mcq{206}{Simple index number is calculated for:}{One variable}{Two variable}{Three variable}{None}
\mcq{207}{An index number computed for single variable commidity called:}{Aggregative index}{Composite index}{Simple index}{None of these}
\mcq{208}{Simple index numbers involves commodities:}{More than one}{One}{Two}{Three}
\mcq{209}{An index number is called simple index number when it is computed form:}{Simple index}{Single variable}{Two variable}{None}
\mcq{210}{The formula used for simple aggregative method is:}{$\dfrac{\sum P_{n}}{\sum P_{o}}\times100$}{$\dfrac{P_{n}}{P_{o}}\times100$}{$\dfrac{P_{n}}{P_{n-1}}\times100$}{$\dfrac{P_{n}}{P_{n-1}}\times100$}
\mcq{211}{An index number that can serve many purpose:}{General purpose index}{Special purpose index}{Cost of living index}{None of these}
\mcq{212}{Purchasing power of money can be assessed through:}{Simple index}{Fisher's ideal index}{Consumer price index}{Volume index}
\mcq{213}{The most suitable average for index number is:}{A,M}{G.M}{Median}{Mode}
\mcq{214}{Price of base year is denoted by:}{$P_{on}$}{$P_{o}$}{$P_{n}$}{$Q_{o}$}
\mcq{215}{$\dfrac{\sum p_{n}q_{n}}{\sum p_{o}q_{o}}\times100$ is equal to:}{$P_{on}$}{$I_{n}$}{$Q_{on}$}{$V_{on}$}
\mcq{216}{Laspeyre's indec number:}{$\dfrac{\sum p_{o}q_{o}}{\sum p_{n}q_{o}}\times100$}{$\dfrac{\sum p_{n}q_{o}}{\sum p_{o}q_{o}}\times100$}{$\dfrac{\sum p_{n}q_{n}}{\sum p_{o}q_{o}}\times100$}{$\dfrac{\sum p_{o}q_{n}}{\sum p_{o}q_{o}}\times100$}
\mcq{217}{Laspeyre's formula is also called:}{Base year weighted}{Ideal}{Current year weighted}{None of these}
\mcq{218}{Base year weighted index no. are also called:}{Laspeyre}{passches}{Fisher}{None of these}
\mcq{219}{Which of the following is called an ideal index no. :}{Laspeyre}{Passches}{Fisher}{None of these}
\mcq{220}{If all the values considered in calculating an index are of equal importance, then index is:}{Weighted}{Simple}{Unweighted}{None of these}
\mcq{221}{Passche's formula is also called:}{Ideal index}{Base year weighted}{Current year weighted}{None of these}
\mcq{222}{If lespeyre price index number is 200, Passche's price index number is 200, then Fisher price index number is:}{100}{200}{40000}{None of these}
\mcq{223}{Fisher index number is of Laspeyre's and passches andex number is:}{A.M}{G.M}{Median}{Mode}
\mcq{224}{Current year weighted index no. are also called:}{Laspeyre}{Passches}{Fisher}{Marshal}
\mcq{225}{G.M of Laspeyre's and passche's index number:}{Fisher}{Marshall}{Both}{None}
\mcq{226}{Fisher's ideal index written as:}{$\sqrt{\dfrac{P}{L}}$}{$\sqrt{\dfrac{L}{P}}$}{$\sqrt{L\times P}$}{None of these}
\mcq{227}{If laspeyre's index = 116 and Passche's index = 110. Then Fisher's index is equal to:}{110}{116}{113}{112.96}
\mcq{228}{CPI stands for:}{Consumer price index}{Chain price index}{Complete price index}{None of these}
\mcq{229}{The first index was calculated in:}{Italy}{America}{Germany}{England}
\mcq{230}{The price of wheat as a single commodity are compared by:}{Simple index}{Composite index}{Price relatives}{Link relatives}
\mcq{231}{An ordinary arrangement of object is called:}{Combination}{Permutation}{Sample space}{Factorial}
\mcq{232}{A non-ordinary arrangement of things is called:}{Permutation}{Combination}{Order pair}{Set}
\mcq{233}{${}^{n}P_{r}=\ldots\ldots$ :}{$\dfrac{n!}{r!}$}{$\dfrac{n!}{(n+r)!}$}{$\dfrac{n!}{r!(n-r)!}$}{$\dfrac{n!}{(n-r)!}$}
\mcq{234}{${}^{n}C_{r}=\ldots\ldots$:}{${}^{n-1}C_{r}$}{${}^{n}C_{n-r}$}{${}^{n}C_{n-r-1}$}{None of these}
\mcq{235}{${}^{n}C_{r}=\ldots\ldots$:}{$\dfrac{n!}{(n-r)!}$}{$\dfrac{(n-r)!}{n!}$}{$\dfrac{r!(n-r)!}{n!}$}{$\dfrac{n!}{(n-r)!\,r!}$}
\mcq{236}{${}^{n}C_{o}=\ldots\ldots$:}{0}{1}{n!}{n}
\mcq{237}{${}^{8}C_{3}=\ldots\ldots$:}{20}{56}{6720}{None of these}
\mcq{238}{${}^{6}C_{3}=\ldots\ldots$:}{6}{15}{20}{25}
\mcq{239}{$\dbinom{10}{3}=$:}{10}{3}{120}{100}
\mcq{240}{If A=\{1,2,3\} then power set of A contain:}{6 subsets}{7 subsets}{5 subsets}{9 subsets}
\mcq{241}{A set which conatin countless elements is called:}{Sub set}{Universal set}{Finite set}{Infinite set}
\mcq{242}{IF $A=\{3,4,5$, $B=\{5,6,7\}$ the $A\cap B$ is:}{$\{3,4,5,6,7\}$}{$\{5,6,7\}$}{$\{3,4\}$}{$\{5\}$}
\mcq{243}{${}^{3}P_{2}$ is equal to:}{3}{5}{6}{1}
\mcq{244}{$A\cup A=\ldots\ldots$:}{A}{S}{$A\subseteq B$}{S-A}
\mcq{245}{If A=\{points on a line\} then A is called:}{Finite set}{Infinite set}{Null set}{Power set}
\mcq{246}{If A is the subset of B then we may write:}{$A\subseteq B$}{$A\supset B$}{$A\subset B$}{$A\in B$}
\mcq{247}{If A=\{3,4,5\}, B=\{5,6,7\} then $(A\cap B)$ is:}{\{3,4,5,6,7\}}{\{6,7\}}{\{3,4\}}{\{5\}}
\mcq{248}{If A=\{a,b,c\}, B=\{d,e,f\} then A-B is:}{A}{B}{\{e,f\}}{None of these}
\mcq{249}{The set of all subset a set is called:}{Universal set}{Empty set}{Power set}{Subset}
\mcq{250}{If A=\{1,2,3\} then power set of A contains:}{6 subsets}{7 subsets}{8 subsets}{9 subsets}
\mcq{251}{A set, which conatins no element is called:}{Subset}{Sample space}{Null set}{None of these}
\mcq{252}{If \{0\} then the number of sebsets of A are:}{1}{2}{3}{Zero}
\mcq{253}{A' is defined as:}{S-A}{A-S}{S+A}{None of these}
\mcq{254}{If A=\{2,3\}, B=\{5,6\} then A x B contains:}{2 elements}{3 elements}{4 elements}{5 elements}
\mcq{255}{If ``n'' is the number of elements of a set then the total numbr of subsets of the this set is:}{2n}{$n^{2}$}{$2^{n}$}{n!}
\mcq{256}{If ``n'' is the number of elements of a set then the total numbr of subsets of the this set is:}{2n}{$n^{2}$}{$2^{n}$}{n!}
\mcq{257}{$A\cup A'=\ldots\ldots$:}{A}{S}{S-A}{$\phi$}
\mcq{258}{A unit set is also called:}{Universal set}{Subset}{Empty set}{Singleton set}
\mcq{259}{$A\cap A'=\ldots\ldots$ :}{A}{S}{S-A}{$\phi$}
\mcq{260}{$(A\cup B)'=\ldots\ldots$ :}{$A'\cup B'$}{$(A\cap B)'$}{$A'\cap B'$}{None of these}
\mcq{261}{$A\cup(B\cap C)=\ldots\ldots$:}{$(A\cap B)\cup(A\cap C)$}{$(A\cap B)\cap(A\cup C)$}{$(A\cup B)\cup(A\cup C)$}{$(A\cup B)\cap(A\cup C)$}
\mcq{262}{$\dbinom{10}{3}=\ldots\ldots$:}{10}{3}{120}{100}
\mcq{263}{The number of permutations of all the ``n'' distinct objects are:}{\{n-1\}!}{n!}{\{n+1\}!}{n}
\mcq{264}{${}^{n}P_{r}=\ldots\ldots$ :}{$\dfrac{n!}{r!}$}{$\dfrac{n!}{(n-r)!}$}{$\dfrac{n!}{(n+r)!}$}{$\dfrac{n!}{(n-r)!\,r!}$}
\mcq{265}{Four books in 1st, 2nd, 3rd and 4th position can be arranged in:}{4 ways}{6 ways}{12 ways}{24 ways}
\mcq{266}{The number of permutations of ``n'' distinet objects arranged in a circle is:}{[n+1]!}{[n-1]!}{n!}{$n^{2}$}
\mcq{267}{5 persons can sit around a table:}{20 ways}{24 ways}{25 ways}{30 ways}
\mcq{268}{The possible arrangements of word STATISTICS are:}{$\dfrac{10!}{3!3!2!1!1!}$}{$\dfrac{10!}{4!2!1!}$}{$\dfrac{10!}{2!2!2!}$}{None of these}
\mcq{269}{${}^{n}C_{r}=\ldots\ldots$:}{$\dfrac{n!}{r!}$}{$\dfrac{n!}{(n-r)!}$}{$\dfrac{n!}{(n+r)!}$}{$\dfrac{n!}{(n-r)!\,r!}$}
\mcq{270}{The relationship between ${}^{n}C_{r}$ and ${}^{n}P_{r}$ is:}{${}^{n}P_{r}={}^{n}C_{r}$}{${}^{n}P_{r}=\dfrac{1}{r}\times{}^{n}C_{r}$}{${}^{n}P_{r}=r!\times{}^{n}C_{r}$}{None of these}
\mcq{271}{If B=$\{x\mid x$ is prime factor of $3\}$ then B must be:}{\{3\}}{\{0\}}{$\phi$}{Subset}
\mcq{272}{$A'\cup(B\cap C)=\ldots\ldots$:}{$(A'\cap B)\cup(A'\cap C)$}{$(A'\cap B)\cap(A'\cup B)$}{$(A'\cup B)\cap(A'\cup C)$}{$(A'\cup B)\cap(A'\cap B)$}
\mcq{273}{$(A\cap B)'=\ldots\ldots$:}{$A'\cup B'$}{$(A\cup B)'$}{$A'\cap B'$}{None of these}
\mcq{274}{Combination of two basic experiments is called:}{Simple experiment}{Random experiment}{Compound experiment}{Scientific experiment}
\mcq{275}{Two dice are rolled together, the total possible resilts are:}{$6^{2}$}{$2^{6}$}{12}{62}
\mcq{276}{If three coins are tossed then np. of elements in a set S are:}{3}{8}{9}{6}
\mcq{277}{The number of term in the expansion of the binomial $(a+b)^{n}$ is:}{n}{n-1}{n+1}{2n}
\mcq{278}{The possible arrangements of word ASSASINATION are:}{3326400}{332640}{33264}{33264000}
\mcq{279}{The number of ways in which 11 players can be selected from 15 players is:}{${}^{15}C_{11}$}{${}^{15}P_{11}$}{${}^{15}C_{15}$}{${}^{11}C_{11}$}
\mcq{280}{$A\cup A'=\ldots\ldots$}{A}{S}{A-S}{$\phi$}
\mcq{281}{$A\cap S=\ldots\ldots$:}{A}{S}{S-A}{$\phi$}
\mcq{282}{$A\cap(A\cup C)=\ldots\ldots$:}{$(A\cup B)\cap(B\cup C)$}{$(A\cap B)\cup(A\cup C)$}{$A\cup B\cup C$}{$A\cap B\cup C$}
\mcq{283}{The set of natural numbers is a:}{Universal set}{Finite set}{Infinite set}{Exact set}
\mcq{284}{When we say two sets are disjoint, it means:}{They have all common elements}{They have no common elements}{Their union produces universal set}{Their intersection produces universal set}
\mcq{285}{If we have a set A=\{1,2,3,4\} then power set of A contains:}{4 elements}{14 elements}{16 elements}{24 elements}
\mcq{286}{$\dbinom{10}{3}+\dbinom{10}{4}=\ldots\ldots$ :}{${}^{11}C_{4}$}{${}^{12}C_{4}$}{${}^{20}C_{7}$}{${}^{11}C_{3}$}
\mcq{287}{${}^{n}C_{r}+{}^{n}C_{r+1}=\ldots\ldots$:}{${}^{n}C_{r-1}$}{${}^{n}C_{r+2}$}{${}^{n+1}C_{r}$}{${}^{n+1}C_{r+1}$}
\mcq{288}{The number of permutations of 5 distinct objects are:}{[5-1]!}{5!}{[5+1]!}{5}
\mcq{289}{${}^{6}P_{2}=\ldots\ldots$:}{$\dfrac{6!}{2!}$}{$\dfrac{6!}{(6-2)!}$}{$\dfrac{6!}{(6+2)!}$}{$\dfrac{6!}{(6+2)!\,6!}$}
\mcq{290}{The number of permutations of 6 distinet objects arranged in a circle are:}{[6+1]!}{[6-1]!}{6!}{$6^{2}$}
\mcq{291}{${}^{7}C_{2}=\ldots\ldots$:}{$\dfrac{7!}{2!}$}{$\dfrac{7!}{(7-2)!}$}{$\dfrac{7!}{(7+2)!}$}{$\dfrac{7!}{(7+2)!\,2!}$}
\mcq{292}{The numbers of terms in the expansion of $(a+b)^{2}$ is:}{5}{4}{6}{7}
\mcq{293}{Three coins are tossed together, the total possible results are:}{$3^{2}$}{$3^{3}$}{$2^{3}$}{$3^{4}$}
\mcq{294}{If a coin and a die is rolled together the total possible results are:}{12}{16}{8}{6}
\mcq{295}{Three books of different colours are to be arranged in a book-self. The possible arrangements are:}{3}{6}{1}{2}
\mcq{296}{The word Random can be arranged in:}{24 ways}{120 ways}{720 ways}{6 ways}
\mcq{297}{n(n-1)(n-2) \ldots\ldots 3.2.2 is equal to:}{$\sum n$}{$\sum n(n-1)$}{n!}{n}
\mcq{298}{${}^{n}C_{r}=\ldots\ldots$:}{${}^{n-1}C_{r}$}{${}^{n}C_{n-r-1}$}{${}^{n-1}C_{r}$}{None of these}
\mcq{299}{${}^{n}P_{r}=\ldots\ldots$:}{$\dfrac{n!}{r!}$}{$\dfrac{n!}{(n+r)!}$}{$\dfrac{n!}{r!(n-r)!}$}{$\dfrac{n!}{(n-r)!}$}
\mcq{300}{If three coins are tossed then the possible outcomes are:}{3}{9}{4}{8}
\mcq{301}{The probability of an event always liew between:}{0 and 1}{-1 and 1}{-1 and zero}{None of these}
\mcq{302}{The idea of probability was given by:}{Cricket game}{Basket ball game}{Chess game}{Gambling game}
\mcq{303}{The probability of all possible outcomes of a random experiment is always equal to:}{Zero}{One}{Infinity}{All of these}
\mcq{304}{Which is an impossible event when a die is rolled:}{2 or 3}{5 or 6}{1}{9 or 7}
\mcq{305}{The probability of an impossible event is equal to zero:}{Always}{Never}{Oficenty}{None of these}
\mcq{306}{Probability is the measurement of:}{Certainty}{Chances}{Uncertainty}{Both B and C}
\mcq{307}{Probability of a sure event is:}{Zero}{$\dfrac{1}{2}$}{$\dfrac{3}{4}$}{1}
\mcq{308}{The probability of an event can never be:}{Positive}{Less then one}{Negative}{None}
\mcq{309}{A particular result of an experiment is called:}{Trail}{Simple event}{Compound event}{Outcome}
\mcq{310}{The probability of an event must be:}{Negative}{Non negative}{Greater than 1}{Less then zero}
\mcq{311}{If A and B are two independent events then:}{P[A/B]= P[A]}{P[A] = P[B]}{P [B/A] = P[B]}{BothA and C}
\mcq{312}{A coin and dice thrown together in:}{Two ways}{12 ways}{8 ways}{None of these}
\mcq{313}{If $P(A)=\dfrac{1}{3},P(B)=\dfrac{2}{3}$ if A and B are independent then P$P(A\cap B)$ is:}{$\dfrac{2}{6}$}{$\dfrac{2}{9}$}{0}{None of these}
\mcq{314}{If $P(A\cap B)\neq0$ then events A and B are:}{Not mutually exclusive event}{Mutually exclusive event}{Equally likely events}{Independent events}
\mcq{315}{The probability of appearing a tail / head, when a fair coin is tossed:}{0}{$\dfrac{1}{2}$}{$\dfrac{1}{4}$}{1}
\mcq{316}{If two events have not common sample points, they are called:}{Mutually exclusive}{Equally likely}{Simple event}{Compound event}
\mcq{317}{If A and B are mutually exclusive events then:}{$A\cap B=\phi$}{$A\cap B\neq\phi$}{$A\cup B=\phi$}{$A\cup B\neq\phi$}
\mcq{318}{Any sub set of the sample space is called:}{Element}{Sample point}{Impossible event}{Event}
\mcq{319}{Compound event contains outcome:}{One}{Less than}{More than One}{Zero}
\mcq{320}{If $P(A\cap B)\neq o$ then A and B are events:}{Simple}{Equally likely}{Independent}{Not mutually exclusive}
\mcq{321}{In throwing three coins, possible outcomes are:}{2}{4}{16}{8}
\mcq{322}{When two dice are rolled the possible outcomes are:}{6}{12}{24}{36}
\mcq{323}{The probability of getting a black card is:}{$\dfrac{13}{52}$}{$\dfrac{4}{52}$}{$\dfrac{26}{52}$}{$\dfrac{12}{52}$}
\mcq{324}{If $A\cap B\neq\phi$ then A and B are:}{Disjoint}{Mutually Exclusive}{Not mutually}{Both A and B}
\mcq{325}{Probability of occurring even number on a die is:}{$\dfrac{1}{2}$}{$\dfrac{1}{3}$}{$\dfrac{1}{6}$}{None of these}
\mcq{326}{If P(A) = P(B), then A and B are events:}{Mutually exclusive}{Independent}{Dependent}{Equally likely}
\mcq{327}{If three cards are drawn from a pack of 52 cards the sample space will be:}{${}^{52}C_{3}$}{${}^{52}C_{0}$}{${}^{3}C_{0}$}{None of these}
\mcq{328}{Probability of an impossible event is:}{One}{Two}{Zero}{$\dfrac{1}{2}$}
\mcq{329}{If a coin is tossed twice, then the probability of getting one head and one tail is:}{$\dfrac{1}{2}$}{$\dfrac{1}{6}$}{$\dfrac{1}{4}$}{$\dfrac{2}{3}$}
\mcq{330}{If n coins are tossed, the possible outcomes:}{n}{2}{$2^{n}$}{All of these}
\mcq{331}{Two Dice are rolled together, the total possible results are:}{$6^{2}$}{$2^{6}$}{12}{62}
\mcq{332}{In tossing of two perfect coins the probability at least one head is:}{0}{$\dfrac{1}{4}$}{$\dfrac{3}{4}$}{1}
\mcq{333}{If P(A) = 0.7 then $P(\overline{A})$ is:}{0.70}{0}{1}{0.30}
\mcq{334}{For two independent events A and B if P(A) = 0.12, P(B) = 0.2, then $P(A\cap B)$:}{0.32}{0.06}{0.024}{0.24}
\mcq{335}{The set of all possible outcomes of a random experiment is called:}{Null set}{Samples space}{Empty}{Impossible}
\mcq{336}{The probability of drawing any one spade from pack of 52 cards is:}{$\dfrac{1}{52}$}{$\dfrac{1}{13}$}{$\dfrac{4}{13}$}{$\dfrac{1}{4}$}
\mcq{337}{If $P(A\cap B)=0.25,\ P(A)=0.75$ then $P\left[\dfrac{B}{A}\right]$:}{1}{Zero}{1/3}{2/3}
\mcq{338}{If A and B are mutually exclusive events then $P(A\cap B)$ is:}{1}{0.5}{Zero}{None of these}
\mcq{339}{If three coins are tossed then the probability of at least one head is:}{$\dfrac{7}{8}$}{$\dfrac{2}{8}$}{$\dfrac{3}{8}$}{Zero}
\mcq{340}{The total number of possible outcomes when two coins are tossed are equal to:}{3}{4}{6}{8}
\mcq{341}{The probability of getting head on a single throw of a coin is:}{$\dfrac{1}{2}$}{$\dfrac{1}{4}$}{$\dfrac{1}{3}$}{None of these}
\mcq{342}{IF two dice are rolled, then P(same number)=:}{$\dfrac{1}{36}$}{$\dfrac{1}{6}$}{$\dfrac{1}{9}$}{$\dfrac{3}{4}$}
\mcq{343}{If a single dice is thrown, total possiblr outcomes are:}{36}{6}{1}{0}
\mcq{344}{The probability of drawing a picture card from a pack of 52 cards is:}{$\dfrac{12}{52}$}{$\dfrac{26}{52}$}{$\dfrac{4}{52}$}{$\dfrac{10}{52}$}
\mcq{345}{If a die is rolled once, the probability of add number is equal to:}{$\dfrac{1}{2}$}{$\dfrac{1}{3}$}{$\dfrac{1}{4}$}{$\dfrac{1}{6}$}
\mcq{346}{The probability of drawing red card out of 52 cards is:}{$\dfrac{1}{2}$}{$\dfrac{13}{52}$}{$\dfrac{4}{52}$}{Zero}
\mcq{347}{The probability of an event can not be:}{-0}{$<1$}{$>1$}{-1}

\end{document}
